\begin{definition}[Lector de firma y transporte de representantes]
\label{def:firma555}
Sea \(\mathcal X_\times=D_9^2\times\{N,E,S,O\}\), de cardinal \(324\).
El avance \(P\) desplaza una casilla en la dirección almacenada y la media
vuelta \(H\) invierte esa dirección. El lector utiliza
\[
\ell_9(1,2,4,8,7,5)=(0,1,2,3,4,5),\qquad
\ell_9(3)=\ell_9(6)=\ell_9(9)=0.
\]
En cada una de seis ventanas de nueve pasos se evalúa
\(\ell_9(\rho_9(ij))\) antes del avance producto seleccionado por TRIT en
los tiempos \(2,5,8\pmod9\). La suma de cada ventana se reduce módulo
\(10\). Tras los pasos \(27\) y \(54\), la dirección se invierte.
La regla determina \(E_{\rm sig}(x)\) a partir de la posición y dirección
de \(x\), con las convenciones de tiempo del núcleo.
\end{definition}

\begin{proposition}[Refinamiento mínimo estable de la firma]
\label{prop:firma555}
Partiendo de la equivalencia por firma \(\sim_0\), la iteración
\[
x\sim_{n+1}y\ \Longleftrightarrow\
x\sim_n y,\quad Px\sim_n Py,\quad Hx\sim_n Hy
\]
termina y produce la congruencia estable más gruesa que refina la firma.
En el dominio anterior el censo es \(26\to28\to28\); la fibra de
\(555555\) se separa en tres clases de ocho representantes.
\end{proposition}
\begin{proof}
En cada paso se refina una partición de un conjunto finito. El proceso se
estabiliza tras un número finito de escisiones. En el punto fijo, la
equivalencia se preserva bajo \(P\) y \(H\). Si una congruencia estable
refina \(\sim_0\), por inducción refina cada \(\sim_n\): sus imágenes
por ambos operadores permanecen equivalentes. Refina, por tanto, el punto
fijo obtenido, lo que prueba su minimalidad en información añadida.

La enumeración finita aplica la definición a las \(81\) posiciones y las
cuatro orientaciones. Se agrupan primero sus seis sumas módulo \(10\), y
después la terna formada por la clase de \(x\), la de \(Px\) y la de
\(Hx\), hasta estabilizar. El histograma resultante es
\(27\cdot8+108=324\), y la única escisión de firma es la indicada.
Los siguientes representantes hacen explícitas las tres imágenes:
\[
\begin{array}{c|c|c}
x&E_{\rm sig}(x)&E_{\rm sig}(Px)\\\hline
(3,5,N)&555555&555177\\
(3,4,N)&555555&555717\\
(3,4,S)&555555&555771
\end{array}
\]
La regla completa y su enumeración ejecutable se conservan en el paquete
de reproducción; los testigos de la tabla muestran además que conservar
solamente la firma no hace unívoca su continuación.
\end{proof}

La conclusión describe la información mínima del transporte, no una
nueva masa. El factor aditivo conserva dieciocho representantes por firma
(nueve posiciones compatibles y dos orientaciones). Al combinarlo con las
fibras producto de tamaños \(24\) y \(8\), se obtienen respectivamente
\(432\) y \(144\). La incidencia regional y la memoria de hoja conservan
así funciones distintas dentro del mismo catálogo.
